\section{Runtime and Empirical Results}
\label{sec:prof-results}

The decisive empirical question for ORIUS is not whether one domain achieves the lowest
cost or the smallest interval width. It is whether the runtime layer materially reduces
true-state violations under degraded observation while remaining operationally cheap and
auditable. The current repository truth supports that claim most strongly in the battery
witness row and then at bounded depth across AV, industrial, and healthcare.

The battery witness establishes the strongest result surface: ORIUS repair-aware
controllers eliminate true-state violations under the locked degraded-telemetry protocol
while maintaining bounded intervention rates. The cross-domain runtime evidence then
shows that the same runtime grammar remains cheap enough to be plausible across
non-energy domains. The important result is therefore joint: safety repair is doing real
work, and the added runtime layer is not computationally extravagant.

\begin{table*}[t]
\centering
\scriptsize
\caption{Runtime and governance summary for the professor-facing ORIUS draft.}
\label{tab:prof-runtime-governance}
\begin{tabularx}{\textwidth}{@{}p{0.24\textwidth}p{0.16\textwidth}p{0.14\textwidth}p{0.14\textwidth}p{0.12\textwidth}p{0.14\textwidth}@{}}
\toprule
\textbf{Domain} & \textbf{Row} & \textbf{P95 latency (ms)} & \textbf{Fallback (\%)} & \textbf{Audit / lifecycle (\%)} & \textbf{Boundary} \\
\midrule
Battery Energy Storage & Witness row & 0.812 & 100.0 & 100.0 & deepest theorem-to-artifact closure\\
Autonomous Vehicles & Defended bounded row & 0.744 & 100.0 & 100.0 & bounded TTC plus predictive-entry-barrier contract\\
Medical and Healthcare Monitoring & Defended bounded row & 0.701 & 100.0 & 100.0 & bounded monitoring and alert-release semantics\\
\bottomrule
\end{tabularx}
\end{table*}


Three results matter most.
\begin{enumerate}[leftmargin=1.25em,itemsep=0.2em]
  \item \textbf{Safety reduction:} the battery witness row is the clearest proof that
  reliability-conditioned repair prevents hidden violations that remain visible under
  deterministic and static-interval baselines.
  \item \textbf{Runtime plausibility:} the defended rows remain in the sub-millisecond
  regime at the P95 level, which means ORIUS behaves like a lightweight safety layer
  rather than a second controller.
  \item \textbf{Governed release:} certificates, fallback paths, and audit continuity
  are observed rather than assumed, so runtime safety is measured as a released-action
  discipline instead of an undocumented filter.
\end{enumerate}

The combined runtime-governance table also makes the evidence boundary easy to read.
Battery is the witness row. AV, industrial, and healthcare are defended bounded rows
with full certificate rates and bounded runtime semantics. Navigation and aerospace are
present in the runtime surface, but they remain blocked from defended closure for data
and task reasons rather than architecture reasons. That is exactly the disciplined claim
shape a professor-facing draft should preserve.

What this means scientifically is that ORIUS is being evaluated as a release layer, not
as a task policy. The battery witness demonstrates the deepest theorem-to-runtime
closure, but the runtime-governance table shows the more general point: intervention,
fallback, latency, and certificate validity can be measured under one contract across
qualitatively different physical systems. That is why the paper uses bounded rows and a
witness row rather than forcing every domain into one artificial leaderboard.

\begin{table*}[htbp]
\centering
\scriptsize
\caption{Deployment validation scope for the active three-domain ORIUS draft.}
\label{tab:orius-deployment-validation-scope}
\resizebox{\textwidth}{!}{%
\begin{tabular}{p{2.5cm}p{2.5cm}p{2.2cm}p{2.2cm}p{2.8cm}p{3.0cm}}
\toprule
\textbf{Deployment surface} & \textbf{Governing artifact} & \textbf{Scope type} & \textbf{Current status} & \textbf{Bounded manuscript claim} & \textbf{Exact non-claim or gap}\\
\midrule
Battery witness runtime & reports/publication/dc3s\_main\_table.csv & witness\_replay & bounded\_reference & Battery supports the deepest runtime and theorem witness surface. & This is still not unrestricted field deployment.\\
Autonomous-vehicle defended replay & reports/orius\_av/full\_corpus/runtime\_summary.csv & bounded\_replay & defended\_bounded & AV supports bounded replay under the TTC plus predictive-entry-barrier contract. & It does not claim full autonomous-driving closure.\\
Healthcare defended replay & reports/healthcare/runtime\_governance\_summary.csv & bounded\_replay & defended\_bounded & Healthcare supports bounded monitoring and alert-release claims. & It does not claim regulated clinical deployment and still uses governance-pass proxy wording for certificate-valid release.\\
OOD and adversarial completeness & chapters/ch34\_outside\_current\_evidence.tex; reports/publication/adversarial\_probing\_robustness\_table.csv & explicit\_non\_claim\_register & bounded\_non\_claim & The monograph can discuss bounded active probing and non-claim discipline. & It does not claim universal adversarial completeness or unrestricted OOD safety.\\
\bottomrule
\end{tabular}}
\end{table*}


The deployment-scope table prevents the results section from sliding into unbounded
deployment rhetoric. Even the strongest battery row stops short of unrestricted field
deployment. The defended bounded rows stop at replay-backed and runtime-governed
validation. The blocked rows remain visible precisely so that the paper's universal
architecture claim stays larger than a single benchmark number but smaller than a
universal deployment claim.

One further interpretive point matters for the professor review draft. ORIUS is not
claiming that every domain should optimize the same surface metric. Some rows care about
dispatch legality, others about headway or envelope preservation, others about
intervention legality. The universal result is instead that each domain can expose a
true-state safety object, and that ORIUS can reduce risk on that object by governing
release under degraded observation. That is the empirical thesis the paper can defend
without overstating parity.
